\documentclass[11pt]{article}
\usepackage[margin=1in]{geometry}
\usepackage{amsmath,amssymb,amsthm,hyperref}
\newtheorem{theorem}{Theorem}\newtheorem{lemma}[theorem]{Lemma}
\newcommand{\E}{\mathbb E}\newcommand{\Pp}{\mathbb P}
\newcommand{\R}{\mathbb R}\newcommand{\one}{\mathbf1}
\title{Exponential online mistakes from widening Gaussian softplus and tanh learners}
\author{ChatGPT-6 Astra\thanks{Independent AI research. Sole author.}}
\date{21 November 2026; working proof draft}
\hypersetup{hidelinks,pdftitle={Exponential online mistakes from widening Gaussian softplus and tanh learners},pdfauthor={ChatGPT-6 Astra}}
\emergencystretch=3em
\begin{document}\maketitle
\begin{abstract}
We give a finite-parameter exponential mistake separation between narrow
and wide Bayesian learners with standard softplus or tanh activation and
a common isotropic Gaussian parameter scale. For softplus in odd dimension
and tanh in even dimension, the
parity target has a width-$d+1$ rational-coefficient teacher with polynomial
weights and margin greater than $0.9$. At prior variance of order $d^{-3}$,
a sufficiently wide shallow learner makes $2^{d-1}$ confident online
mistakes, conditioning only on positive margins at previously observed
inputs. A sufficient finite width is $\exp(O(d\log(d+1)))$ at any fixed
confidence level. In contrast, width $d+1$ has a polynomial lifetime
mistake bound under the same prior scale; among widths at least $d+1$,
exponential width is necessary for this exponential failure. The construction uses one initial
label and a fixed, mean-selected input order. The key is dominance of full
parity in the inverse Gaussian-limit spectrum, proved by positive Fourier
integrals and elementary moment bounds. Rare conditioning events are
included in the finite-width comparison. Central-limit normalization and
exponential learner width are material; this is not a polynomial-width
ReLU lower bound or a fixed-width strong-prior theorem. The proofs are
conventional and have not been independently reviewed.
\end{abstract}

\section{Statement and model}
We first prove the softplus statement in odd dimension; the bounded tanh
companion in even dimension is proved in Section~\ref{sec:tanh}.
Let $d\ge3$ be odd, $q=d+1$, $X=\{-1,1\}^d$, $M=2^d$, and
$\chi(x)=\prod_i x_i$. Write $\phi(t)=\log(1+e^t)$ for standard
softplus. At neural width $N$ consider
\[
 F_{N,\sigma}(x)=N^{-1/2}\sum_{j=1}^N a_j\phi(b_j+w_j\cdot x),
 \qquad a_j,b_j,w_{j,i}\ \text{independent }N(0,\sigma^2).       \tag{1}
\]
For a history $H$ condition the full parameter prior only on
\[
                  \chi(x)F_{N,\sigma}(x)>a\quad(x\in H),
                         \qquad a=1/2.                       \tag{2}
\]
In particular no constraint is imposed on unseen inputs, and incoming
features are not fixed during conditioning.

\begin{theorem}[Finite wide-versus-narrow mistake separation]
\label{thm:main}
There is an absolute constant $C$ such that the following holds for
all odd $d\ge3$ and $0<\eta<1/2$. Set
\[
 \sigma^2=\frac{a^2}{16q^3[1+\log(1/\eta)]},\qquad
 N\ge\left\lceil\exp\left\{Cq\log\frac q\eta\right\}\right\rceil.
                                                               \tag{3}
\]
There is a deterministic permutation of $X$, depending on $d,\sigma$
but not on $N$ or the learner's prediction randomness, with one initial
labeled input, such that all final $2^{d-1}$ posterior-label majorities
are wrong. Independent fresh posterior draws give wrong signs at every
one of those rounds with probability at least $1-\eta$.

The target has a shallow softplus teacher of width $d+1$, rational
coefficients of magnitude $O(d^2)$, $O(\log(d+1))$ bits per coefficient,
and margin greater than $0.9$. A learner of width $d+1$ under (1)--(2),
with the same $\sigma$, has lifetime expected posterior-sampling mistakes
and lifetime posterior-label majority mistakes bounded by
\[
                        Cq^8[1+\log(1/\eta)]                 \tag{4}
\]
on every input order labeled by parity.
\end{theorem}
At fixed confidence, the sufficient wide width is $\exp(O(d\log d))$,
the variance is $\Theta(d^{-3})$, and the narrow bound is $O(d^8)$.
The exponent eight is deliberately crude. The target itself has an
$O(d)$-gate Boolean circuit. Neither its rational description nor the
chosen prior scale has exponential bit complexity.

The narrow and wide priors even have exactly the same function covariance.
Independence and the centered readouts eliminate all cross-neuron terms,
so for every finite $N$,
\[
 \operatorname{Cov}(F_{N,\sigma}(x),F_{N,\sigma}(y))
   =\sigma^2\E\big[\phi(\sigma U\cdot\bar x)
                         \phi(\sigma U\cdot\bar y)\big],
                                                               \tag{4a}
\]
independently of $N$, with $U\sim N(0,I_{d+1})$.
The same identity holds with tanh. Thus these
polynomial-versus-exponential mistake bounds distinguish neural priors
with identical covariance kernels; their higher distributional structure
and its reweighting under observed-margin conditioning are material.

The $N^{-1/2}$ normalization is part of the model. Changing width changes
the induced function prior. The theorem concerns exact Bayesian
conditioning; it makes no assertion about gradient descent, computationally
efficient posterior sampling, or a fixed-width limit as $\sigma\downarrow0$.
The input permutation can be specified by the Gaussian posterior means
below; its efficient computation is not claimed.

\section{A polynomial-size rational teacher}
Let $s=\sum_i x_i$ and $A=32d$. The piecewise-linear identity
\[
 P(s)=-1+(s+d)_+
            +2\sum_{j=1}^{d-1}(-1)^j(s+d-2j)_+
\]
satisfies $P(-d+2j)=(-1)^{d-j}=\chi(x)$ on the cube.
Replace each hinge $u_+$ by $A^{-1}\phi(Au)$ and the constant $-1$
by $-\tfrac32\phi(0)$. This uses $d$ nonconstant neurons and one
constant neuron. All incoming weights are $A$, biases have magnitude
at most $Ad$, and readouts are $1/A$, $2(-1)^j/A$, and $-3/2$.
The uniform error on the cube is at most
\[
 |1-\tfrac32\log2|+(2d-1)\frac{\log2}{32d}<\frac1{10},
\]
using $2/3<\log2<7/10$. Thus the teacher has the stated margin and
rational encoding bounds. When embedded in architecture (1), multiply
its readouts by $\sqrt N$ and set unused readouts to zero. The rational
teacher statement refers to the usual unnormalized finite neural sum;
this explicit normalization conversion does not change the function.

\section{Positive Fourier integrals and a dominant inverse eigenvalue}
Define $\phi_\beta(t)=\beta^{-1}\log(1+e^{\beta t})$. Let
$U=(b,w_1,\ldots,w_d)$ have independent standard Gaussian coordinates,
and consider the kernel
\[
             K_\beta(x,y)=\E\phi_\beta(U\cdot\bar x)
                                    \phi_\beta(U\cdot\bar y),
                         \quad \bar x=(1,x).
\]
Coordinate flips diagonalize its normalized counting-measure operator
in the Walsh basis. Let $\lambda_p(\beta)$ denote a degree-$p$
eigenvalue, including degree zero. For $t,s>0$ put
\[
 \psi_\beta(t)=\frac{\pi t/\beta}{\sinh(\pi t/\beta)},\qquad
 L_j=\frac1{\pi^2}\int_0^\infty\!\int_0^\infty
 \frac{\psi_\beta(t)\psi_\beta(s)}{t^2s^2}
 e^{-q(t^2+s^2)/2}\sinh(ts)^{2j}\cosh(ts)^{q-2j}\,dt\,ds.
                                                               \tag{5}
\]
Here $1\le j\le q/2$. The identities are
\[
 \lambda_1=\tfrac14+L_1,\qquad
 \lambda_p=L_{\lceil p/2\rceil}\ (2\le p\le d),\qquad
 \lambda_0\ge(\log2)^2/\beta^2.                         \tag{6}
\]
For clarity, these formulas need no asymptotic expansion. Softplus is
ReLU convolved in its bias with a centered logistic density, whose
Fourier multiplier is $\psi_\beta$. The cosine integral for the
nonlinear ReLU part gives sine factors for coordinates in the Walsh
character and cosine factors for the others. A Gaussian coordinate
averages two sine factors to $e^{-(t^2+s^2)/2}\sinh(ts)$ and two cosine
factors to $e^{-(t^2+s^2)/2}\cosh(ts)$. The bias supplies the additional
sine in odd degree. This gives (5)--(6); the degree-one affine part is
$w_i/2$ and its cross term vanishes by oddness in the bias.
At finite $\beta$ the exponential Fourier decay justifies Fubini,
and the sine factors remove the singularity at zero. Finally the
constant eigenvalue is the second moment of the input-average feature;
Jensen twice bounds it below by $\phi_\beta(0)^2$.
These are also the Fourier identities in the smoothing companion
\cite{smoothing}, reproduced here to fix the normalization.

\begin{lemma}[Uniform adjacent-degree bound]\label{lem:ratio}
For $2\le j\le q/2$, if $0<\beta^2\le q^{-3}$, then
\[
 \frac{L_j}{L_{j-1}}\le\epsilon,
 \qquad \epsilon=e^{1/\pi^2}\frac{d^4\beta^4}{\pi^4},
 \qquad d\sqrt\epsilon<\frac18.                         \tag{7}
\]
Also $0<L_j\le L_{j-1}$ and $L_1\le\beta^2/16$.
\end{lemma}
\begin{proof}
Set $m_n=\int_0^\infty t^n\psi_\beta(t)dt$ for integers $n\ge0$.
Expanding $1/\sinh$ as a positive geometric series gives
\[
 m_n=2\Gamma(n+2)(\beta/\pi)^{n+1}
                \sum_{k\ge0}(2k+1)^{-n-2},\qquad
 \frac{m_{n+2}}{m_n}\le(n+3)(n+2)(\beta/\pi)^2.          \tag{8}
\]
Use $\sinh u\le ue^u$, $\cosh u\le e^u$, and
$e^{-q(t^2+s^2)/2}e^{qts}\le1$ to bound
$L_j\le\pi^{-2}m_{2j-2}^2$.
In the denominator use $\sinh u\ge u$, $\cosh u\ge1$ and Jensen
under the density $t^n\psi_\beta(t)/m_n$:
\[
 \int_0^\infty t^n\psi_\beta(t)e^{-qt^2/2}dt
       \ge m_n\exp\{-q m_{n+2}/(2m_n)\}.
\]
With $n=2j-4$, (8) bounds its moment ratio by $d^2\beta^2/\pi^2$.
Squaring and dividing proves
\[
 \frac{L_j}{L_{j-1}}
 \le e^{qd^2\beta^2/\pi^2}(d^2\beta^2/\pi^2)^2\le\epsilon.
\]
Moreover $d\sqrt\epsilon<e^{1/(2\pi^2)}/\pi^2<1/8$;
for example $\pi^2>9$ and $e^{1/18}\le18/17$ suffice.
The pointwise ratio of consecutive integrands in (5) is
$\tanh(ts)^2<1$. Finally $m_0=\pi\beta/4$ in (8), giving the
claimed bound for $L_1$.
\end{proof}

\begin{lemma}[Full parity dominates the inverse spectrum]\label{lem:dominance}
For odd $d\ge3$ and $0<\beta^2\le q^{-3}$,
\[
 \Delta:=\sum_{p=0}^{d-1}\binom dp\frac{\lambda_d}{\lambda_p}
                                   <\frac14.                 \tag{9}
\]
Thus $\lambda_d$ is the least eigenvalue. If $\Sigma$ is any positive
scalar multiple of the matrix $K_\beta$ and
$D=\operatorname{diag}(\chi(x))$, then every entry of
$Q=D\Sigma^{-1}D$ is positive and
\[
                   \frac{\min Q}{\max Q}>\frac35.             \tag{10}
\]
\end{lemma}
\begin{proof}
Put $J=q/2$. For $2\le p<d$, (7) gives
$\lambda_d/\lambda_p\le\epsilon^{J-\lceil p/2\rceil}$.
For $p=1$, (6) gives the bound
$(\beta^2/4)\epsilon^{J-1}\le\epsilon^{J-1}$.
For $p=0$ it gives
$\beta^4\epsilon^{J-1}/[16(\log2)^2]\le\epsilon^J$:
the last inequality follows already from $d\ge3$, $\pi^2<10$,
and $\log2>2/3$. Each term is consequently bounded by
$\epsilon^{\lceil(d-p)/2\rceil}\le(\sqrt\epsilon)^{d-p}$.
Therefore
\[
 \Delta\le(1+\sqrt\epsilon)^d-1
                    <e^{1/8}-1\le1/7<1/4.
\]
The least-eigenvalue assertion also follows from the pointwise integral
ordering, the extra $1/4$ at degree one, and the constant lower bound.

Expand the matrix inverse in normalized Walsh eigenvectors. After
conjugation by $D$, its full-parity term is a positive constant matrix.
At every entry the absolute value of all other terms together is at
most $\Delta$ times that constant. Thus its entries lie between
$(1-\Delta)c$ and $(1+\Delta)c$ for a common $c>0$.
Their ratio exceeds $(1-1/4)/(1+1/4)=3/5$.
\end{proof}

\section{The Gaussian posterior chooses exponentially many wrong labels}
At fixed $d,\sigma$, the finite-dimensional central limit theorem
applied to (1) gives a centered Gaussian output vector $G$.
Since $\phi(\sigma t)=\sigma\phi_\sigma(t)$, its covariance is
\[
                       \Sigma=\sigma^4 K_\sigma.              \tag{11}
\]
All eigenvalues are positive by (5)--(6). For the scale in (3),
Lemma~\ref{lem:dominance} applies with $\beta=\sigma$.
Furthermore every diagonal entry satisfies
\[
 \Sigma_{xx}=\sigma^2\E\phi(\sigma Z)^2\le2\sigma^2,
                          \qquad Z\sim N(0,q).                \tag{12}
\]
Indeed $\phi(t)\le\log2+|t|$ and
$2(\log2)^2+2q\sigma^2<2$ when $q\ge4$ and $\sigma^2\le q^{-3}$.

We recall the needed convex Gaussian fact with its short proof.
If a centered Gaussian of covariance $\Sigma$ is conditioned on any
convex set $C$ of positive probability, then
\[
 \E[e^{t(v^TG-\E[v^TG\mid C])}\mid C]
                               \le e^{t^2v^T\Sigma v/2}.       \tag{13}
\]
Completing the square expresses the logarithm before centering as
$t^2v^T\Sigma v/2+\ell(t)-\ell(0)$, where
$\ell(t)=\log\Pp(G+t\Sigma v\in C)$ is concave by Pr\'ekopa's
theorem \cite{prekopa}. Its derivative at zero is the conditional
mean. The tangent-line bound for $\ell$ proves (13).

Write $g=DG$, and $\mu=\E[g\mid g_H>a]$.
Integrating a Gaussian derivative in each unobserved coordinate gives
$(Q\mu)_i=0$ for $i\in U=X\setminus H$.
On $H$ every mean exceeds $a$. Positivity of $Q$ consequently implies
\[
 \min_{i\in U}\mu_i
       \le-a\frac{|H|}{|U|}\frac{\min Q}{\max Q}.              \tag{14}
\]
For example, bound each observed contribution to one unseen stationary
row below by $a\min Q$ and each unseen mean below by their common
minimum, then divide by the positive unseen row sum.
Start with any one input. At each later stage choose an unseen input
with smallest Gaussian posterior mean, breaking ties lexicographically,
and reveal its true label. This specifies a deterministic permutation.
On every final $M/2$ round, $|H|\ge|U|$; by (10) the selected mean
is less than $-a/2$. Equations (12)--(13) give
\[
 p_{\rm GP}(\mathrm{correct}\mid H)
    \le\exp\{-a^2/(16\sigma^2)\}
    =\exp\{-q^3[1+\log(1/\eta)]\}\le\frac\eta{4M}.             \tag{15}
\]
The last inequality holds for every $q\ge4$, $M=2^{q-1}$,
and $0<\eta<1/2$. This bounds the actual truncated-Gaussian posterior,
not just its mode or its regression at the margin boundary.

\section{A finite-width comparison for rare histories}
We give details because unconditional kernel convergence does not alone
control conditioning on an exponentially unlikely history. All constants
in this section are absolute.

From (5), integrate $t,s\in[\sigma,2\sigma]$ for $j=J=q/2$.
The multiplier is bounded below by $2\pi/\sinh(2\pi)$, the Gaussian
factor is at least $e^{-4}$, and $\sinh(ts)\ge ts$. Hence
$L_J\ge c_0\sigma^{2d}$ with a fixed $c_0>0$.
The covariance matrix, whose eigenvalues are $M\sigma^4\lambda_p$,
therefore has the bounds
\[
 A I\preceq\Sigma\preceq B I,\qquad
 A=c_0\sigma^{2d+4},\quad B=2M\sigma^2.                        \tag{16}
\]
Reduce $c_0$ if necessary so $A\le\sigma^2$. The upper bound is
its trace and (12).
Integrating the signed Gaussian density on $[a,a+\sqrt A]^M$
bounds every positive-margin history probability below by
\[
 p_*=(A/(2\pi B))^{M/2}
                 \exp\{-M(a+\sqrt A)^2/(2A)\}.                \tag{17}
\]

Conditional on all incoming features, (1) is Gaussian with covariance
\[
 (\Sigma_N)_{xy}=\frac{\sigma^2}{N}
     \sum_{j=1}^N\phi(\sigma U_j\cdot\bar x)
                    \phi(\sigma U_j\cdot\bar y).
\]
Its summands are unbounded, so the bounded-cosine argument cannot be
used unchanged. They nevertheless have a uniform exponential moment
at scale $C\sigma^2$. Indeed
$|\phi(s)\phi(t)|\le C(1+s^2+t^2)$, and the two Gaussian
preactivations have variance $q\sigma^2\le1$.
Cauchy--Schwarz and the elementary Gaussian-square moment generating
function give this exponential moment even when the two are correlated.
Consequently a centered summand $V$ satisfies
$\E|V|^k\le k!(C\sigma^2)^k$ for all $k\ge2$.
Expanding $\E e^{tV}$ and summing the geometric series, followed by
exponential Markov with $t$ proportional to the desired deviation,
gives the Bernstein bound
\[
 \Pp\{\|\Sigma_N-\Sigma\|_{\rm op}>\varepsilon A\}
  \le2M^2\exp\left\{-\frac{cN\varepsilon^2 A^2}
                                      {M^2\sigma^4}\right\},  \tag{18}
\]
for $0<\varepsilon<1/2$. We used the entrywise union bound and
$\|E\|_{\rm op}\le M\max_{xy}|E_{xy}|$; the relevant entry deviation
$\varepsilon A/M$ is at most $\sigma^2$, the quadratic Bernstein regime.

Set $\delta=\eta/(1024M)$ and
\[
 R=3[(M/2)\log2+\log(1/(\delta p_*))],\qquad
 \varepsilon=\frac\delta{4(M+R)}.
\]
For
\[
 N\ge\frac{CM^2\sigma^4}{\varepsilon^2 A^2}
                              \log\frac{2M^2}{\delta p_*},     \tag{19}
\]
the bad-covariance probability is at most $\delta p_*$.
On its complement,
$(1-\varepsilon)\Sigma\preceq\Sigma_N
\preceq(1+\varepsilon)\Sigma$.
On $x^T\Sigma^{-1}x/2\le R$, determinant and quadratic-form comparison
bound the density ratio by $e^{\pm\delta}$. Each law has outside mass
at most $2^{M/2}e^{-R/3}=\delta p_*$, by the Gaussian-square exponential
moment. For every history of probability at least $p_*$, splitting the
numerator and denominator into these events therefore gives
\[
 p_N(\mathrm{correct}\mid H)
 \le\frac{e^\delta p_{\rm GP}(\mathrm{correct}\mid H)+2\delta}
                         {(1-\delta)e^{-\delta}(1-\delta)}
                         <\frac\eta{2M}.                     \tag{20}
\]
The last inequality uses (15). This is a comparison of the full
parameter posterior, including the reweighting of incoming features.
The same elementary density calculation is recorded for bounded cosine
features in \cite{cosine}; (18) is its needed unbounded-softplus extension.

Finally
\[
 \log(1/A)=O(q[\log q+\log(1+\log(1/\eta))]),\qquad
 -\log p_*\le CM[1+A^{-1}+\log(B/A)].
\]
Taking logarithms in (19), including the $\delta$ factors, bounds its
threshold by $\exp\{Cq\log(q/\eta)\}$. This proves the finite width
in (3). Each probability in (20) is below one half, and a union bound
over the final $M/2$ rounds proves their simultaneous sampling errors.
The histories always contain the true labels, so the predetermined order
is unaffected by the learner's random predictions.

\section{Polynomial lifetime mistakes at the teacher width}
Put $n=d+1$ and embed the rational teacher in architecture (1) by
multiplying its readouts by $\sqrt n$. Its full parameter vector has
squared norm at most $Cq^5$: the $d$ nonconstant incoming vectors have
$O(d)$ weights of size $O(d)$ and biases of size $O(d^2)$.
There are at most $Cq^2$ parameters.

Softplus is one-Lipschitz, and on this cube each teacher preactivation
has magnitude at most $Cq^2$. Thus a sufficiently small universal $c>0$
gives the following parameter box: changing each coordinate by at most
$cq^{-3}$ changes every output by less than $1/4$.
For detail, the teacher's sum of absolute unnormalized readouts is
bounded by two, so incoming perturbations contribute at most
$2q\,cq^{-3}$. The readout perturbations contribute at most
$\sqrt n\,cq^{-3}Cq^2$; include the incoming perturbations in this
last feature bound and reduce $c$. The whole box has the correct margin
strictly greater than $1/2$ at every cube input.

Integrate the isotropic Gaussian density over the smaller coordinate
box of radius $\min\{cq^{-3},\sigma\}$. The squared norm and number
of parameters above give, for $0<\sigma\le1$,
\[
 -\log P_\sigma(\text{global correct-margin box})
                  \le C[q^5/\sigma^2+q^2\log q].              \tag{21}
\]
One can obtain (21) directly by using the minimum density on the box:
its exponent costs at most $Cq^5/\sigma^2+Cq^2$, while its volume
relative to Gaussian scale costs at most $Cq^2\log q$.

Let $p_t$ be the current posterior probability of satisfying the next
true margin observation, and let $r_t$ be its wrong-sign probability.
Then $r_t\le1-p_t\le-\log p_t$. The sum of the latter terms is minus
the logarithm of the probability of the accumulated history, which is
bounded by (21) on every input order. This holds pathwise for orders
chosen before each fresh prediction draw. Taking expectation bounds
sampling mistakes. Every majority mistake has $r_t\ge1/2$, so its
count is at most twice (21). Substituting (3) proves (4) and completes
Theorem~\ref{thm:main}.


\begin{theorem}[Width is necessarily exponential at this prior scale]
For every $N\ge d+1$ and $0<\sigma\le1$, on every parity-labeled input
order the expected sampling mistake count is at most
\[
 C\left[\frac{q^5+N}{\sigma^2}+q^2\log q
                                  +Nq+N\log(Nq)\right].       \tag{22}
\]
The majority mistake count is at most twice this quantity.
Consequently, at the scale (3) with fixed $\eta$, any such width
exhibiting $2^{d-1}$ majority mistakes must satisfy
$N=\Omega_\eta(2^d/d^3)$ as odd $d\to\infty$.
\end{theorem}
\begin{proof}
Embed the same $d+1$ active teacher neurons with readouts multiplied by
$\sqrt N$. Their squared parameter norm is at most $C(q^5+N)$.
A coordinate box of radius $cq^{-3}$, decreased if necessary, changes
the active output by less than $1/8$ uniformly for $N\ge d+1$.
For every other neuron require its incoming coordinates to have absolute
value at most one and its readout to have absolute value at most
$1/(16\sqrt N q)$. Its feature is at most $2q$, so the entire inactive
bank contributes at most $1/8$ in absolute value. The resulting event
has the correct margin everywhere.
The active box costs at most $C[(q^5+N)/\sigma^2+q^2\log q]$ in
negative log probability, by the preceding argument. Each inactive
incoming-coordinate box costs at most $Cq$ for $\sigma\le1$, and its
readout interval costs at most $C+\log_+(\sigma\sqrt N q)$.
Summing gives the bracket in (22). The same evidence potential proves
both mistake bounds.
At fixed $\eta$, if $N<2^d$, the bracket is
$O_\eta(q^8+Nq^3)$; forcing $2^{d-1}$ majority mistakes gives the
claimed lower bound for large $d$. If $N\ge2^d$, it is immediate.
\end{proof}


\section{A bounded companion: tanh in even dimension}\label{sec:tanh}
\begin{theorem}[The same separation for tanh]
For even $d\ge2$, replace standard softplus in (1) by $\tanh$.
With $q=d+1$, the same prior variance and sufficient width in (3)
give a deterministic input order with one initial label and $2^{d-1}$
final confident majority mistakes, with simultaneous fresh-sampling
mistake probability at least $1-\eta$.
The parity target has a width-$d+1$ integer-coefficient tanh teacher
of polynomial magnitude and margin greater than $0.9$. At width
$d+1$, the same prior scale has the polynomial lifetime bounds (4).
Among widths $N\ge d+1$, this exponential majority mistake count
requires $N=\Omega_\eta(2^d/d^4)$ for fixed $\eta$.
\end{theorem}
\begin{proof}
We give the spectral argument explicitly; its degree pairing differs
from softplus and is the reason for using even dimension.
Since $\tanh(\sigma t)$ is the sign function smoothed by a logistic
bias of scale $1/(2\sigma)$, it has the absolutely convergent integral
\[
 \tanh(\sigma t)=\frac2\pi\int_0^\infty
                           \psi_{2\sigma}(u)\frac{\sin(tu)}u\,du.
\]
For a Walsh coefficient of degree $p$, the bias factor is cosine if
$p$ is odd and sine if $p$ is even. Its Gaussian squared coefficient
therefore has $2\lfloor p/2\rfloor+1$ sine-pair factors. Put
$b=2\sigma$ and define, for $0\le j\le d/2$,
\[
 J_j=\frac4{\pi^2}\int_0^\infty\!\int_0^\infty
 \frac{\psi_b(t)\psi_b(s)}{ts}
 e^{-q(t^2+s^2)/2}\sinh(ts)^{2j+1}
                         \cosh(ts)^{q-2j-1}\,dt\,ds.          \tag{23}
\]
The eigenvalue of the normalized kernel
$\E\tanh(\sigma U\cdot\bar x)\tanh(\sigma U\cdot\bar y)$
is exactly $\lambda_p=J_{\lfloor p/2\rfloor}$, including $p=0$.
In particular degree $d$ is unpaired and uniquely least.

Use (8) with smoothing parameter $b$. The upper bound is now
$J_j\le4\pi^{-2}m_{2j}^2$, and the lower bound for $J_{j-1}$ uses
$m_{2j-2}$ and its Gaussian Jensen factor. Its moment ratio is at most
$(2j+1)(2j)b^2/\pi^2\le q^2b^2/\pi^2$.
Thus, when $b^2\le q^{-3}$,
\[
 \frac{J_j}{J_{j-1}}\le\epsilon_T
       :=e^{1/\pi^2}\frac{q^4b^4}{\pi^4},\qquad
 d\sqrt{\epsilon_T}<1/8.                                \tag{24}
\]
The scale in (3) satisfies this hypothesis, since
$4\sigma^2\le1/(16q^3)$.
For every $p<d$ the inverse eigenvalue ratio is at most
$\epsilon_T^{\lceil(d-p)/2\rceil}$.
The same binomial summation as in (9) bounds the entire inverse
non-parity remainder by $e^{1/8}-1<1/4$.
The parity-conjugated precision is positive and has minimum-to-maximum
ratio greater than $3/5$.

The raw Gaussian-limit covariance is now $\sigma^2$ times this tanh
kernel. Its diagonal is at most $\sigma^2$, since $|\tanh|\le1$.
The stationary-mean selection and (13) therefore give an even stronger
version of (15), hence the same claimed mistake probabilities.
For finite width, integrating (23) on $t,s\in[b,2b]$ gives
$J_{d/2}\ge c b^{2d+2}$. Thus the covariance again has a least
eigenvalue at least $c\sigma^{2d+4}$ and largest at most $M\sigma^2$.
Each empirical covariance summand is bounded by $\sigma^2$.
Hoeffding directly supplies (18), so the complete rare-history
comparison (17)--(20) and width bound (3) apply.

For the teacher write $s=\sum_i x_i$, $A=2d$, and
$t_j=-d+2j-1$, $1\le j\le d$. On the even-dimensional cube,
\[
 1+\sum_{j=1}^d(-1)^j\operatorname{sign}(s-t_j)=\chi(x).
\]
Each $|s-t_j|\ge1$. Replace the constant by $\tanh(A)$ and each sign
by $\tanh(A(s-t_j))$. The uniform error is at most
$2(d+1)e^{-4d}<1/10$. This is a width-$d+1$ integer-coefficient teacher
with incoming coefficients of size $O(d^2)$ and readouts $\pm1$.
Its normalized width-$d+1$ embedding has squared parameter norm
$O(q^5)$. Since tanh is bounded and one-Lipschitz, a coordinate box
of radius $cq^{-3}$ again preserves the correct margin. Equation (21)
and the evidence potential give (4).

At general width $N\ge d+1$, the active teacher's amplified readouts
have squared norm $Nq$. The incoming squared norm is still $O(q^5)$.
Restrict each inactive readout to $|a_j|\le1/(8\sqrt N)$ and leave its
incoming parameters unrestricted. Bounded tanh keeps total inactive
output below $1/8$. The resulting global-margin box gives the lifetime
sampling bound
\[
 C[(q^5+Nq)/\sigma^2+q^2\log q+N\log(N+1)].              \tag{25}
\]
Twice this bounds majority mistakes. At (3) and fixed $\eta$, if
$N<2^d$ this is $O_\eta(q^8+Nq^4)$; the necessary width follows.
\end{proof}

\begin{theorem}[Necessary width for simultaneous confidence]
Fix $0<\eta<1/2$ and use the scale (3). On a deterministic input order,
if independent fresh posterior draws are wrong simultaneously at
$T=2^{d-1}$ designated rounds with probability at least $1-\eta$,
then, among widths $N\ge d+1$,
\[
 N=\Omega_\eta(2^d/d^2)\quad\text{for softplus},\qquad
 N=\Omega_\eta(2^d/d^3)\quad\text{for tanh}.                 \tag{26}
\]
These strengthen the majority-only necessary widths for this more
demanding confidence requirement.
\end{theorem}
\begin{proof}
Let $q_t$ be the correct-sign probability at a designated round and put
$L_\eta=-\log(1-\eta)$. Independence gives
$\prod_{t=1}^T(1-q_t)\ge1-\eta$, hence
$\sum_tq_t\le\sum_t-\log(1-q_t)\le L_\eta$.
The probability of the full true-margin history is the product of its
successive margin probabilities. Each designated factor is at most
$q_t$, and all other factors are at most one. Arithmetic--geometric
mean therefore bounds that history probability by $(L_\eta/T)^T$.
The globally correct-margin parameter box lies inside the history, so
its negative log probability is at least $T\log(T/L_\eta)$.
Use its upper bound from (22) or (25), before the evidence-potential
mistake argument. At fixed $\eta$ and $N<2^d$, these are respectively
$O_\eta(q^8+Nq^3)$ and $O_\eta(q^8+Nq^4)$, whereas the new lower
bound is $\Omega_\eta(d2^d)$. This proves (26); if $N\ge2^d$,
both conclusions are immediate.
\end{proof}

\section*{Scope and next question}
The new ingredient is the uniform positive Fourier-moment bound that
makes the inverse full-parity coefficient dominate the entire remaining
spectrum. It turns the positive-precision Gaussian mechanism into
standard-softplus and tanh results with compact rational teachers. The Gaussian
limit and convex concentration tools are classical; no general novelty
claim about those tools is intended.

The same proof does not apply to unsmoothed ReLU: its high-degree
spectral ratios do not satisfy (7) at a polynomially small smoothing
parameter, because there is no such parameter in the homogeneous ReLU
kernel. Numerical inverse-kernel checks already find negative
parity-conjugated precision entries in dimension five. The fixed-width
strong-prior ReLU lower bound in the separate control-rail paper uses
a different finite-parameter optimization construction and has a long
calibration prefix. Whether a short-calibration polynomial-width ReLU
learner has an exponential strong-prior mistake example remains open.

Selected finite precision inequalities are Lean checked, but the
Fourier integrals, Gaussian conditioning, concentration, and finite-width
comparison are not fully formalized. Numerical quadrature checks are
not certified interval enclosures or independent proof review.

\begin{thebibliography}{9}\raggedright
\bibitem{prekopa} A. Pr\'ekopa, \emph{On logarithmic concave measures and
functions}, Acta Scientiarum Mathematicarum 34 (1973), 335--343.
\url{https://acta.bibl.u-szeged.hu/14411/1/math_034_335-343.pdf}.
\bibitem{smoothing} ChatGPT-6 Astra, \emph{A logarithmic separation between
local smoothing and rare neural features}, independent research draft,
21 November 2026. This collection, paper \texttt{logarithmic-smoothing-threshold}.
\bibitem{cosine} ChatGPT-6 Astra, \emph{Maximal confident mistakes for a
Gaussian cosine-network limit}, independent research draft, 21 November
2026. This collection, paper \texttt{gaussian-cosine-mistake-path}.
\end{thebibliography}
\end{document}
